\subsection{Primera iteración de la Fase 4: Análisis de los cambios necesarios}
\label{subsec:F4I1}

\subsubsection{Diferentes velocidades por plataforma}

El reloj del Director encadena sus \textit{ticks} con \texttt{setTimeout} y
una pausa fija de 500 milisegundos escrita en el código
(\texttt{PAUSA = 500}). Este valor es bastante correcto para el simulador en
PC\footnote{En mi ordenador, usar 480 es más aproximado al G\&W real.},
pero no tiene por qué serlo para el ESP32: la duración real de cada
\textit{tick} (barrido, motor de estados y renderizado) depende del
procesador, y replicar la velocidad subjetiva de la máquina original puede
exigir pausas distintas en cada plataforma. Además, un valor fijo en el
código obliga a recompilar para probar ritmos distintos, y la fase siguiente
(el generador del cartucho compacto) necesita que el ritmo sea una propiedad
del cartucho para que el sistema sobre ESP32 lo reciba ya resuelto.

\subsubsection{Aceleración del juego con el tiempo}

En las máquinas reales la partida no mantiene un ritmo constante: a medida
que avanza el tiempo, el juego se vuelve más difícil. El simulador actual
mantiene la misma pausa durante toda la partida, de modo que la dificultad
solo crece por acumulación de desenlaces, no por velocidad. Falta el
mecanismo que reduzca la pausa del reloj con el tiempo, de forma que la
partida acelere sola, sin necesidad de modos de juego separados.

\subsubsection{Política de lanzamiento de rutas}

El nacimiento de rutas es hoy una política fija del motor de estados: solo
nace una ruta nueva en el \textit{tick} en que ha finalizado la vida de la
anterior (las animaciones pueden ser simultáneas, pero no son parte del juego
ni son rutas que se lancen). En las máquinas reales la densidad de tráfico es
parte del diseño de cada juego y crece con la partida: comienzan cruzando
pocas rutas y, a medida que avanza, van aumentando las rutas simultáneas.
Falta, por tanto, que esta política se pueda redefinir y suba al cartucho,
editable desde el diseñador gráfico.

\subsubsection{Revisión de campos en el cartucho}
Como hay que añadir todos los temas de configuración general en el cartucho para su uso por el director del juego y por el diseñador de cartuchos compactos, se revisará que incluya toda la información necesaria para generar el cartucho final.
